% /root/monodromy, 2026-11-19. Proved; independent root, obstructions_crystals and reduction_descent full audits passed.
\subsection{Eventual exact groups on every ordinary principally polarized surface}
\label{subsec:ordinary-ppav-surface-full-groups}
\status{Proved and checked. The parent and both other research agents
independently audited the complete assembly, including every exceptional jet
stratum, the exact invariant quotient, the norm degeneration, actual natural
block extraction and the arithmetic image. The cited source scopes were
checked independently; the level-three finite calculation was also rerun.}

Throughout this subsection $(A,\Theta)$ is an ordinary principally polarized
abelian surface in characteristic two, $L=\mathcal O_A(3\Theta)$, and the
linear systems are complete unless ambient hypersurfaces are explicitly
mentioned. A symmetric representative of the polarization is chosen over the
algebraic closure. Translation of that representative gives an isomorphic
geometric family and does not change its geometric group.

\begin{proposition}[All invariant covectors on an ordinary surface]
\label{prop:ordinary-ppav-all-covectors}
For every $d\ge1$ the geometric generic relative rank is $r_d=9d^2-1$.
The entire geometric invariant-covector space has dimension zero for odd $d$
and three for even $d$. In even degree it is canonically the quotient
\[
 \mathbf F_2^{H}/\mathbf F_2\mathbf 1,
 \qquad H=\operatorname{Hom}(\ker F_A,\mathbf G_m),\quad |H|=4,
\]
with its natural Galois action. Equivalently its basis can be indexed by the
three nonzero characters, which Galois permutes.
\end{proposition}
\begin{proof}
The ordinary nodal degeneration of Lemma~\ref{lem:ppav-ordinary-members}
proves generic ordinarity and the rank formula in the indecomposable case.
The ordinary product degeneration proves it in the decomposable case.
The coherent relative bundle on $P=\mathbf P H^0(A,L^d)$ is $\Omega_P^1$:
indeed $H^1(A,L^{-d})=0$, $h^2(A,L^{-d})=9d^2$, and
$h^2(A,\mathcal O_A)=1$, and the multiplication map to the last group is
surjective, dual to the inclusion of the defining section.

We give the complete endomorphism reduction, to specify its scope beyond
products. Normalize an invariant two-form by $C(\omega)=\omega$ and put
$W=H^0(A,L^d)$. An invariant covector on the ordinary open extends across
all prime divisors of $P$: its equation in regular frames is
$M^t\lambda=\lambda^{[2]}$, and a largest pole of order $a>0$ would have
order $2a$ on the right and at most $a$ on the left. Normality and local
freeness extend it across codimension two. The dual Euler sequence then
represents it by a linear operator $D\in\End(W)$ modulo scalar identity.
Its Frobenius defect is a constant multiple of $s^{[2]}$; adding a scalar
identity and solving an Artin--Schreier equation removes that defect.
Thus exact lifts are the operators satisfying
\[
                 \beta(s,D(s))=D(s)^{[2]},\qquad s\in W,
\]
modulo $\mathbf F_2\operatorname{id}$, where $\beta$ is relative Cartier
multiplication and $\beta(s,s)=s^{[2]}$.

The eigenvalue and Jordan argument of
Theorem~\ref{thm:product-char2-all-parities-fixed-quotient} applies without
any product assumption: eigenvalues are zero or one, and a Jordan chain
$Dv=w,Dw=0$ would give the nonzero coefficient $\beta(w,w)=w^{[2]}$ on
substituting $v+tw$. Every exact lift is consequently an idempotent, and
its two eigenspaces are orthogonal for $\beta$.

The bundle $\mathcal F=F_{A*}L^d$ is globally generated. For an
indecomposable surface, second-jet surjectivity for $3\Theta$ from
Theorem~\ref{thm:ppav-three-theta-jets} implies surjectivity on a
Frobenius fibre: in regular parameters its algebra is
$k[x,y]/(x^2,y^2)$, a quotient of the second-jet algebra since
$(x,y)^3\subset(x^2,y^2)$. For higher $d$, multiply by a section of
$L^{d-1}$ nonzero at the point. Such a section exists since $L$ is globally
generated. On a decomposable surface the two degree-three elliptic factors
interpolate the two length-two Frobenius fibres, giving the same assertion.
Finite Frobenius duality gives the perfect symmetric Cartier pairing
$\mathcal F\otimes\mathcal F\to L^{d,(2)}$.
Orthogonality therefore makes each of the two eigenspaces preserve the
kernel of every evaluation map. Global generation and perfectness descend
$D$ to a bundle endomorphism of $\mathcal F$, exactly as in the cited
product proof; this reasoning uses only these two properties.

If $d$ is odd, apply \cite[Theorem~1.2(i), printed p.~45; proof
pp.~43--45]{Oda1971} to $F_A$ and $L^d$.
This theorem is stated for abelian varieties of arbitrary dimension and
allows inseparable isogenies. Its hypothesis is that the polarization map
of $L^d$ restricts to an isomorphism of the scheme-theoretic kernel of the
isogeny onto its image. Here this restriction is the restriction of
$3d\lambda_\Theta$ to $\ker F_A$; it is an immersion since $3d$ is odd
and $\lambda_\Theta$ is an isomorphism. Thus
$\End(F_{A*}L^d)=k$. The only exact lifts are zero and identity, and their
quotient gives no invariant covector.

If $d$ is even, write $L^d=F_A^*M$.
The ordinary splitting of $F_{A*}\mathcal O_A$, equivalently the character
decomposition for $\ker F_A\simeq\mu_2^2$, gives
\[
          F_{A*}L^d=\bigoplus_{\chi\in H}M\otimes T_\chi.
\]
The four summands are distinct ample line bundles differing by nontrivial
degree-zero two-torsion line bundles. There are no homomorphisms between
different summands, and every summand has nonzero sections by ample
vanishing and Riemann--Roch. Hence $\End(\mathcal F)=k^4$.
Their Cartier cross pairings vanish, since trace kills nonzero character
weights. Its four projectors are precisely the independent exact lifts:
each projector satisfies the displayed equation, and any idempotent in
$k^4$ is their sum. Modding out the identity gives the stated three-space.
This splitting is \cite[Lemma~4.5, p.~11]{SannaiTanaka2016}, or directly
the character algebra of the ordinary Frobenius torsor.

The construction by projectors is canonical under Galois. The zero
character is fixed, and the three nonzero characters are permuted.
Modulo the sum of all four projectors, those three projectors are a basis.
This proves the final assertion, including its arithmetic description.
\end{proof}

\begin{theorem}[Hasse geometry and the irreducible common kernel]
\label{thm:ordinary-ppav-Hasse-kernel}
For every ordinary principally polarized abelian surface, the full Hasse
divisor in $|L|$ is geometrically reduced with generic Cartier corank one
on each component. For every $d\ge2$ the Hasse divisor in $|L^d|$ is
geometrically irreducible and reduced, with generic stable loss one.
Its geometric relative group is generated by transvections.
If $K_d$ is the common kernel of its entire invariant-covector space,
then $K_d$ is irreducible, and $V_d$ has no invariant vector.
\end{theorem}
\begin{proof}
The decomposable case was proved in
Theorem~\ref{thm:ordinary-products-Hasse-irreducible} and
Corollary~\ref{cor:ordinary-products-irreducible-kernel}.
In the indecomposable case use the genus-two theta divisor and jet
statements of Theorem~\ref{thm:ppav-three-theta-jets}.
Theorem~\ref{thm:ppav-level3-Cartier-symmetric} makes the degree-one Hasse
polynomial reduced with generic corank one. Its components have smooth
generic curves, since the singular-curve discriminant is irreducible and
has ordinary generic member by Lemma~\ref{lem:ppav-ordinary-members}.

We verify all incidence bounds needed to apply the mixed-tacnode argument,
including the rank-four exceptions. Put $e=d-1\ge1$. Let $H_i$ be a
component of the degree-one Hasse divisor, and choose a general smooth
$C\in H_i$. Its Cartier kernel is one-dimensional. The nonzero tangent
functional $dh_1$ belongs to the dual first-jet space at only finitely many
points of $C$: otherwise choose one such $P$ and another outside the
finite exceptional set $R_P$ of the jet theorem. The two first-jet dual
spaces are then disjoint, a contradiction. Choose $P$ also away from the
zeros of the killed differential and from every theta translate
$\Theta-s$, $s\in A[3]$. An integral curve of class $3\Theta$ is not one
of those theta translates, so these exclusions are finite on $C$.

The affine tangent space $\ker dh_1$ to $H_i$ at $C$ surjects onto the
first jets at $P$. The incidence of $D\in|L^e|$ simply tangent to $C$ at
$P$ therefore dominates $|L^e|$: for an infinitesimal change of $D$, move
$P$ to cancel the value change and use these three freely variable first
jets of $C$ to preserve its value and the common tangent line. This is
the three-equation calculation in
Theorem~\ref{thm:product-Hasse-irred}, on the affine cone; projective
scaling imposes no additional condition.

The good contact locus is nonempty. Every $L^e$ inherits the degree-one
second jets at $P$ and the joint first-jet rank bounds at $P,Q$: multiply
by a section of $L^{e-1}$ nonzero at both points, which exists over the
infinite algebraically closed field. Our choice of $P$ eliminates $R_4$,
so all two-point ranks are at least five and are six except at finitely
many $Q$. Within the two imposed contact conditions at $P$, a singularity
at a different $Q$ imposes total rank five generically with $Q$ varying
in dimension two, and total rank at least four at the finitely many
exceptional $Q$. Both are proper loci. Another contact with $C$ imposes
rank four generically with $Q$ varying in dimension one, and rank at least
three at finitely many exceptions, again proper. Second jets exclude
higher contact and a zero normal derivative at $P$. A general member in
the incidence is thus smooth with just this simple tangency. Dominance
then allows it to be ordinary. A single general ordinary $D$ works for
all the finitely many $H_i$.

For this fixed $D$, the contact divisor $\Delta_D\subset|L|$ is
irreducible: its incidence over $P\in D$ is a projective bundle given by
two independent value/tangent conditions. We justify that its generic
point is one simple contact with smooth $C$. For pairs on $D^2$, the
rank-six locus imposes four contact conditions; its point dimension is
two. The exceptional locus $R\cap D^2$ has dimension at most one, since
$R$ is finite over each first point, and imposes at least three conditions
away from $R_4$. The locus $R_4\cap D^2$ is finite: its first point lies
on some $D\cap(\Theta-s)$. It still imposes at least two independent
value conditions. All three counts give codimension at least two in
$|L|$ for two contacts. A singularity of $C$ at a second point imposes
five conditions generically, with a three-dimensional pair space; on
$R$ it imposes at least four with pair dimension at most one, and on
$R_4$ at least three with pair dimension zero. A singularity at the
contact point and order-three contact each impose three conditions with
that point varying on $D$. Thus they too have codimension at least two.
The generic contact is simple and $C$ smooth. The joint tangency incidence
in $(P,C,D)$ is irreducible and dominates both systems by first-jet
surjectivity, so its two ordinary inverse images meet in a dense open.
For general ordinary $D$, the generic $C\in\Delta_D$ is ordinary.

On the multiplication subspace $s_D|L|\subset|L^d|$, the divisorial
nonordinary support is exactly $\bigcup_i H_i\cup\Delta_D$.
This is the union stable-rank formula
\eqref{eq:product-union-stable-rank}. Nonreduced base members have
codimension at least two by the one-node discriminant theorem. A member
containing $D$ is impossible for $e>1$ and a single projective point for
$e=1$. Each $H_i$ meets $\Delta_D$ at the good mixed contact just
constructed. The pointed Cartier lemma and the direct dualizing-form
calculation of Lemma~\ref{lem:mixed-tacnode-Hasse} apply to any ordinary
abelian surface: its only ambient input is bijective Cartier on
$H^1(\mathcal O_A)$ and $H^2(\mathcal O_A)$, and its deformation input is
first-jet surjectivity. Thus the full Hasse divisor is smooth at every
chosen mixed point.

Every global Hasse component restricts to a nonzero positive-degree
polynomial on this multiplication subspace, since the latter contains
an ordinary point. Its restriction contains one of the displayed prime
divisors. Any component containing $H_i$ or $\Delta_D$ contains their
chosen intersection point. Smoothness there forces all these assignments
to be the same component, even if restriction multiplicities occur.
This proves global irreducibility; that same smooth point proves
multiplicity one. An ordinary pair with a single simple contact has
stable loss exactly one by the union formula. Closedness of the iterated
Frobenius rank loci therefore gives generic stable loss one.
Lemma~\ref{lem:char2-defect-one-inertia} now proves transvection generation.

Finally the coherent bundle is $\Omega^1_{\mathbf P^{9d^2-1}}$.
Lemma~\ref{lem:irreducible-hasse-quotients-trivial} implies that every
nonzero proper invariant subspace has trivial quotient. By
Proposition~\ref{prop:ordinary-ppav-all-covectors}, every such subspace
contains $K_d$, making $K_d$ irreducible. Its dimension is at least $32$,
so an invariant line is impossible. This proves $V_d^G=0$ as well.
\end{proof}

\begin{lemma}[A full degree-one open in surface moduli]
\label{lem:ordinary-ppav-GL8-open}
There is a nonempty open in the ordinary moduli stack of principally
polarized abelian surfaces on which the degree-one relative geometric
image for $|3\Theta|$ is $\GL_8(\mathbf F_2)$.
\end{lemma}
\begin{proof}
Work over an algebraically closed field. Use full level five and a fixed
primitive fifth root of unity. The corresponding fine moduli variety is
smooth and geometrically irreducible and carries its universal polarized
abelian scheme, by the exact statement
\cite[Theorem~7.3.11, printed p.~76]{ConradAbelianVarietiesModuli}.
This is an opened author exposition of the moduli theorem; no claim that
its original GIT or compactification proof has been reread is required.
Near the level-five point of the certified product
$E^2$, $E:y^2+xy=x^3+1$, choose a line bundle representing the principal
polarization after an etale neighbourhood. Such representatives form a
smooth torsor under the dual abelian scheme; rigidification at zero
supplies an actual line bundle after this localization. Do not require
this family of representatives to be symmetric: that additional torsor
need not be etale in characteristic two. Each geometric fibre is a
translate of a symmetric representative, so its geometric group is the
same. Take an irreducible smooth neighbourhood $B$ containing the chosen
point.

The cube of this line bundle has a rank-nine direct image, and its
projective bundle carries a nonempty smooth ordinary open $U$ with
geometrically integral fibres. Its rank-eight full Lang frame cover
$Z\to U$ is finite etale. The fibre at the product point is connected by
the certified degree-one $\GL_8$ theorem.
The normalization argument of Lemma~\ref{lem:ordinary-pair-GL8-open}
proves that $Z$ has geometrically connected generic fibre over $B$:
a finite extension defining distinct generic components gives a finite
normalization of the excellent base; on the normal integral projective
bundle open each open-and-closed component of the finite etale cover
surjects, contradicting the connected product fibre. The geometric
generic fibre is smooth, hence normal, so it is also irreducible.
By \cite[Lemma~37.27.5]{Stacks0559}, after shrinking $B$ the frame fibres
are geometrically irreducible. Their connectedness means full $\GL_8$
monodromy. The etale map from $B$ to fine moduli is open, and forgetting
level is open; its image is the asserted nonempty open in the moduli
stack. Translating the representing theta line bundle leaves each
geometric family isomorphic, so this property is intrinsic to the
principally polarized surface.
\end{proof}

\begin{lemma}[Nodal images and an actual natural block]
\label{lem:ordinary-ppav-isogeny-block}
Let $(A',\Theta')$ be an ordinary principally polarized surface with
full degree-one relative image $\GL_8(\mathbf F_2)$, and suppose
$\phi:A'\to A$ is an odd polarized isogeny with
$\phi^*\lambda_\Theta=n\lambda_{\Theta'}$, $n$ odd.
Then $\deg\phi=n^2$, and for every $d\ge n+1$ the degree-$d$
geometric group on $A$ contains a natural $\GL_8(\mathbf F_2)$ on an
actual direct summand, fixing a complementary summand pointwise.
\end{lemma}
\begin{proof}
All work is geometric. Put $L'=\mathcal O(3\Theta')$ and choose symmetric
representatives on these two fixed surfaces. The kernel of $\phi$ is a
constant odd finite group, acting by translations.
A general smooth ordinary $C'\in|L'|$ is preserved by no nonzero kernel
translation: one preserving the line bundle has a nontrivial projective
action in its very ample embedding and hence a proper fixed locus; one
not preserving the line bundle cannot preserve a member. Thus
$C'\to\phi(C')$ is generically one-to-one.

We prove that its image has only ordinary nodes. On an indecomposable
source, for a fixed nonzero kernel element $g$, the two first jets at
$P,P+g$ have rank six except at finitely many $P$. Indeed failure is
contained in
\[
               3P+g\in\Theta',\qquad 3P+2g\in\Theta'.
\]
The two theta translates are distinct because $g\ne0$ and the principal
theta divisor has trivial translation stabilizer. Their intersection is
finite, and $[3]$ is finite. On the rank-six locus, two value conditions
and proportionality of two nonzero gradients impose three independent
conditions. Varying $P$ adds two dimensions, so tangency is generically
absent. The exceptional finite set already imposes two independent value
conditions and is avoided by a general member.

Any three distinct points on an indecomposable $A'$ impose three
independent value conditions on $|3\Theta'|$. To separate $P_3$ from
$P_1,P_2$, choose theta translates $\Theta'_a,\Theta'_b$ through $P_1,P_2$
respectively and not through $P_3$. These choices are dense opens in two
irreducible theta translates. Their addition map dominates $A'$: sums of
two theta curves cover the surface, because any two ample translates
intersect. Thus one can also require the third translate
$\Theta'_{-a-b}$ to avoid $P_3$. The product of these three equations is
a section of $3\Theta'$ vanishing at $P_1,P_2$ but not $P_3$. Cyclically
repeat the construction to prove the three-value rank assertion.
For $P,P+g,P+h$, three independent conditions with $P$ varying in dimension
two exclude triple identifications. There are only finitely many kernel
pairs. On a decomposable source the already proved cubic-product
first-jet and Segre three-value argument gives the identical nodal
conclusion, including its one-dimensional collinear exception. Therefore
$Z=\phi(C')$ is integral nodal with normalization $C'$. For each $g$,
$C'\cap(C'-g)$ has degree $L'^2=18$; its node cover is a quotient of a
union of such degree-eighteen covers. Its permutation group is consequently
a subquotient of a product of copies of $S_{18}$.

The image line bundle is $\operatorname{Nm}_\phi L'$.
Its Neron--Severi class is $3n\Theta$: push forward
$\phi^*\Theta\equiv n\Theta'$ and cancel $n$ in the torsion-free
Neron--Severi group. Both this norm and $L^n$ are symmetric, so their
quotient is a geometric two-torsion line bundle. Since $3n$ is odd,
$\lambda_{L^n}$ identifies the geometric two-torsion of $A$ with that of
its dual. Translating the whole image family by one constant two-torsion
point therefore corrects this quotient and places $Z$ in the exact
linear system $|L^n|$. Denote the corrected finite etale map again by
$\phi$; the translation does not change its kernel identifications.

Fix a general smooth ordinary $D\in|L^{d-n}|$.
Its inverse image $D'=\phi^{-1}D$ is smooth, connected and integral:
it is a smooth ample divisor on the abelian surface, whose connectedness
also follows from ample vanishing in its structure-sheaf sequence.
For varying $C'$, the intersection with $D'$ has degree
\[
                         N=18n(d-n)\ge18.
\]
Its permutation group is the full $S_N$. Here are the required incidence
checks, so that this conclusion does not assume a generic linear-series
monodromy theorem. One-point and ordered two-point incidences are
irreducible projective bundles over $D'$ and $D'^2\setminus\Delta$,
respectively, because $L'$ separates one and two values.
A simple double zero and otherwise simple zeros is attainable. On an
indecomposable source choose $P\in D'$ outside the finite intersections
of $D'$ with the theta translates occurring in the projection of $R_4$.
These intersections are finite since $D'$ is not a theta translate.
The joint surface jets at $P,Q$ have rank at least five for every
$Q\ne P$, and rank six except at finitely many $Q$; their restrictions
to the two tangent directions of $D'$ have rank at least three and
four respectively. After imposing the two conditions for a double zero
at $P$, a second double zero is therefore a proper locus: the generic
second point varies in dimension one and adds two conditions, and each
exceptional point adds at least one. Second jets at $P$ exclude order
three. On a product source choose $P$ also outside the flex fibres;
the same bounds follow from the cubic-product jet lemma, since the
exceptional fibres meet the integral ample $D'$ finitely.
The resulting length-two collision is smooth in the point incidence
and its parameter, since first jets are generated.
The single-collision criterion
\cite[Proposition~3.1, pp.~7--8 of the opened version]{Entin2021}
then gives $S_N$, including characteristic two. Equivalently its proof
uses the irreducible one- and two-point incidences and the nontrivial
transposition inertia of this one ordinary double zero.

After shrinking the coefficient space, $Z$ and $D$ meet transversely,
the image nodes of $Z$ avoid $D$, and $Y=Z\cup D$ is ordinary nodal in
$|L^d|$. The node-avoidance condition is proper: two values at $P,P+g$
with the image point restricted to $D$ impose two conditions with a
one-dimensional point parameter. The normalization sequence and the
isomorphism $H^1_{\et}(A,\mathbf F_2)\xrightarrow{\sim}
H^1_{\et}(A',\mathbf F_2)$ under an odd isogeny give
\[
 0\subset T\subset T_0\subset V_d,\qquad
 T_0/T\text{ constant},\qquad V_d/T_0=Q,
 \qquad\dim Q=8,
\]
where $Q$ is the degree-one relative system of $C'$ and $T$ is graph
cohomology. Its image on $Q$ is full $\GL_8$ even after this shrinking.

The joint node permutation image projects onto $S_N$, with kernel a
subgroup of a product of $S_{18}$'s. The latter has no prime divisor
$31$, whereas $\GL_8(\mathbf F_2)$ does. The factors of $S_N$ are
$A_N$ and cyclic groups of order two; $A_N$ is not $\GL_8(\mathbf F_2)$,
since $11$ divides $|A_N|$ for $N\ge18$ and does not divide
$|\GL_8(\mathbf F_2)|$. Thus the joint image has no composition factor
isomorphic to that simple group. The same is true of its graph image
and of the entire $T_0$ image: the latter only adds a subgroup of the
additive two-group $\Hom(T_0/T,T)$. This step does not assume that the
constant graded pieces split as local systems.

Goursat now shows that the kernel of the entire $T_0$ action still
surjects onto $\GL(Q)$. Inside its block parabolic the translation
kernel is $Q^*\otimes T'$ for a subspace $T'\subset T_0$.
The elementary $H^1(\GL_8(\mathbf F_2),Q^*)=0$ of
Lemma~\ref{lem:char2-group-inputs} splits the remaining cocycle modulo
$T'$ after conjugation. The full translation kernel then removes the
remaining $T'$ columns, producing an actual block-diagonal $\GL_8$.
The whole boundary lies in the coherent ordinary open, so its image is
a subgroup of the same group obtained on the dense smooth ordinary
open. This proves the lemma.
\end{proof}

\begin{theorem}[Eventual exact images for all ordinary principally polarized surfaces]
\label{thm:ordinary-ppav-eventual-full}
Let $(A,\Theta)$ be an ordinary principally polarized abelian surface over
any perfect field $k$ of characteristic two, with a line bundle representing
$\Theta$ over $k$, and embed it by $L=\mathcal O_A(3\Theta)$.
After a finite extension of $k$, there exists an odd polarized isogeny
$\phi:A'\to A$ with multiplier $n$ whose source has degree-one relative
geometric group $\GL_8(\mathbf F_2)$.
For the complete systems $|L^d|$, put $d_1=n+1$.
For the ambient hypersurface systems of this fixed embedding, replace it by
\[
           d_1=\max\{n+1,\operatorname{reg}(\mathcal I_A)-1\}.
\]
For every $d\ge d_1$, put $r=9d^2-1$. If $d$ is odd, the entire geometric
and arithmetic images both equal $\GL_r(\mathbf F_2)$.
If $d$ is even, put $N=r-3$ and let $\Gamma\subset S_3$ be the actual
Galois image permuting the three nonzero characters of $\ker F_A$.
After compatible choices of bases the entire images are
\[
 G_{{\rm geom},d}=
 \left\{\begin{pmatrix}B&T\\0&I_3\end{pmatrix}:
 B\in\GL_N(\mathbf F_2),\ T\in M_{N\times3}(\mathbf F_2)\right\},
\]
\[
 G_{{\rm arith},d}=
 \left\{\begin{pmatrix}B&T\\0&\gamma\end{pmatrix}:
 B\in\GL_N(\mathbf F_2),\ T\in M_{N\times3}(\mathbf F_2),\
 \gamma\in\Gamma\right\},
\]
where $\Gamma$ acts by its three-dimensional permutation matrices.
The representation is irreducible exactly in odd degree.
The integer $n$, and hence the threshold, can be found by a terminating
procedure over finite fields or effectively presented finitely generated
fields of definition; over arbitrary perfect fields the statement gives
existence without presupposing computable coefficients.
\end{theorem}
\begin{proof}
Chai's theorem \cite[Theorem~2, printed p.~477]{Chai1995} states over
any algebraically closed characteristic-$p$ field that the fixed
$\ell$-power polarized Hecke orbit of every ordinary principally
polarized abelian variety is dense in its moduli stack.
Use $g=2$, $p=2$, $\ell=3$. Its orbit meets the nonempty open of
Lemma~\ref{lem:ordinary-ppav-GL8-open}. By the defining form of this orbit
there is an odd polarized isogeny from a source in that open to $A$;
its multiplier is an odd $n$ and its degree is $n^2$. Its finitely many
defining algebraic coefficients lie in a finite extension of $k$.
This auxiliary extension is only used for the geometric subgroup
construction; the arithmetic image below is computed over the original $k$.

For $d\ge n+1\ge2$, Theorem~\ref{thm:ordinary-ppav-Hasse-kernel} gives
an irreducible common kernel $K_d$ and transvection generation.
Lemma~\ref{lem:ordinary-ppav-isogeny-block} places a natural eight-space
$Q$ inside $K_d$, since a natural $\GL_8$ has no nonzero map to a trivial
module. On $K_d$ the same subgroup fixes a complementary summand.
No invariant nondegenerate alternating form can exist: the natural
$\GL_8$ has no invariant bilinear form and no invariant covector, so
both its internal pairing and its pairing with that trivial complement
are zero. Thus $Q$ would lie in the radical. The opened transvection
classification \cite[Theorem~3.6, printed p.~22]{KantorSeress2001} now
forces the image on $K_d$ to be the full $\GL(K_d)$.
For odd $d$ this is already the entire claimed image.

For even $d$, the translation kernel in
$\Hom(V_d/K_d,K_d)$ is $K_d\otimes U$ for a subspace
$U\subset(V_d/K_d)^*$. The natural-module $H^1$ vanishing of
Lemma~\ref{lem:char2-group-inputs} conjugates the remaining extension
to its block form. Its invariant-vector space has dimension
$3-\dim U$. Theorem~\ref{thm:ordinary-ppav-Hasse-kernel} proved that
there is no invariant vector, so $U$ is the whole three-space and
all translations occur. This proves the geometric image.

Proposition~\ref{prop:ordinary-ppav-all-covectors} identifies the entire
geometric covector space together with its Galois action. Consequently
the arithmetic image is contained in the displayed preimage of $\Gamma$.
It contains the full geometric kernel just proved and projects onto
$\Gamma$, because the parameter open is geometrically connected and
its fundamental group surjects onto the absolute Galois group.
It is therefore exactly that preimage. The action on the quotient is
the dual of the covector permutation representation; permutation matrices
identify these two representations. In odd degree the ambient full
linear group leaves no room for a larger arithmetic image.

For the fixed projective embedding, regularity gives
$H^1(\mathbf P^8,\mathcal I_A(d))=0$ when
$d\ge\operatorname{reg}(\mathcal I_A)-1$.
Ambient sections then surject onto $H^0(A,L^d)$; over the nonzero-section
projective open this is an affine bundle, and the pullback Lang cover
has the same geometric and arithmetic image. This avoids a projective
normality assertion in characteristic two.

Finally enumerate all polarized three-power isogenies from the given
surface, by its finite etale three-power torsion and maximal isotropic
subgroups, and reverse them by polarized duals. Compute the degree-one
geometric Lang frame image of each source by the terminating exact
algorithms proved earlier, dovetailing the tests. Hecke density and the
nonempty full-image open ensure that some source passes. Its retained
multiplier is the required $n$. Regularity of the fixed homogeneous ideal
is computable by a finite graded resolution. All field extensions,
finite subgroup schemes and quotients here are described by finite
algebraic data over an effective field of definition. No undecided
Galois-group oracle or unproved generic-group guess enters this procedure.
The odd full module is irreducible, and in even degree its nonzero proper
common kernel of codimension three is invariant, proving the last
representation-theoretic assertion.
\end{proof}

\begin{remark}[Scope of this structural extension]
The theorem treats every ordinary principally polarized abelian surface in
characteristic two, including indecomposable surfaces, after a computable
threshold. It does not identify degree-one groups on every such surface,
and it does not assert a structural full-image formula for arbitrary
ambient smooth projective varieties. Those are separate refinements from
the exact finite algorithms and the general residue theorem already proved.
\end{remark}

\begin{remark}[Finite arithmetic and subgroup checks]
The retained exact program
\path{agents/monodromy/code/ordinary_ppav_group_checks.py}
checks all six $\GL_2(\mathbf F_2)$ character actions on
$\mathbf F_2^4/\mathbf F_2\mathbf1$, their identification with the
three-dimensional permutation representation, and the prime-$31$ and
prime-$11$ separations used for the actual natural block. Its output is
\path{agents/monodromy/outputs/ordinary_ppav_group_checks.json}.
These finite checks support the arithmetic and group calculations; the
geometric incidence arguments and Hecke-density input are proved or cited
separately above.
\end{remark}
